\section{Problem}
\label{section:12-problem}



% 1. Data Size is increasing 
% 2. GraphAI is coming, and doesn't scale with data
% 3. Our current understanding of performance doesn't match the real world (Dynamic Attributed Multigraphs) that we need to operate on under time and resource constraints. 
% 4. We don't have a good understanding of how approximate and incremental algorithms, or heuristic optimizations impact
% 5. We don't have a good method to trade-off between the iron triangle of Scope, Time and Resources, or in this case: Optimality, Time and Cost. 
% 6. There are a lot of acceleration methods and we can't really compare them at the moment. 

\osullikomment{Add motivation here: Tradeoff between time, accuracy and power / resources, and the requirement to understand the relationship between these}

% \nrscomment{youre proposing a hybrid solution- state that}
This thesis proposes a hybrid solution to the problem of \gai benchmarking, where we draw from both existing benchmarking and workload analysis studies to develop an approach that allows for an investigator to trade-off between effectiveness and efficiency when testing a system.
% To advance the adoption and development of \gai, we need to balance the comparability of benchmarks with the insight generated by workload analyses. 
% An ideal benchmark will support the breadth of problems in the \gai taxonomy, and provide the depth of datasets and tasks to effectively characterize their best, worst, expected and typical performance in the real world. 
% It must provide interpretable effectiveness and efficiency metrics for comparing \gai approaches using different abstractions, programming paradigms, system configurations and hardware accelerators. 
A hybrid benchmark must provide sufficient dataset coverage to test the impact of varying input shapes and sizes, support the varying abstractions and representations that drive implementations and collect the metrics required to make informed decisions about the tradeoffs between effectiveness and efficiency. 
It must allow one component to vary while the others remain fixed, and decouple the components sufficiently so that any desired combination of components can be tested. 
